% ============================================================================
% SEC04.TEX - Section 6.4: Physics Material
% ============================================================================

\section{Physics Material}

\begin{learningoutcomes}
\item Membuat Physics Material
\item Mengkonfigurasi friction dan bounce
\item Menerapkan Physics Material ke Collider
\end{learningoutcomes}

\subsection{Properties Physics Material}

\begin{itemize}
\item \textbf{Dynamic Friction}: Friction saat bergerak
\item \textbf{Static Friction}: Friction saat diam
\item \textbf{Bounciness}: Seberapa banyak objek memantul
\item \textbf{Friction Combine}: Cara menggabungkan friction
\item \textbf{Bounce Combine}: Cara menggabungkan bounce
\end{itemize}

\subsection{Contoh Implementasi}

\begin{lstlisting}[language={[Sharp]C}]
// Physics Material dibuat di Project window
// Kemudian diterapkan ke Collider:
public class MaterialExample : MonoBehaviour
{
    public PhysicMaterial bouncyMaterial;
    
    void Start()
    {
        GetComponent<Collider>().material = bouncyMaterial;
    }
}
\end{lstlisting}

\begin{catatan}
Physics Material mempengaruhi bagaimana objek berinteraksi saat collision. Gunakan untuk membuat permukaan licin, kasar, atau bouncy.
\end{catatan}
